% !TEX root = ../../main.tex

% --- KHÔNG ĐƯỢC CÓ LỆNH \documentclass Ở ĐÂY ---

% --- 1. CÁC GÓI CƠ BẢN & ĐỊNH DẠNG ---
\usepackage{amssymb,epsfig,latexsym,array,hhline,fancyhdr} 
\usepackage[normalem]{ulem}
\usepackage{graphicx}
\usepackage{amsmath}
\usepackage{subcaption}
\usepackage{listings}
\usepackage[framemethod=tikz]{mdframed}
\usepackage{float}
\usepackage{tabularx}
\usepackage{tabularray}
\usepackage{vntex}            % Hỗ trợ tiếng Việt
\usepackage{makecell}
\usepackage{changepage}
\usepackage{indentfirst}      % Thụt đầu dòng đoạn văn đầu tiên
\usepackage{algpseudocode}
\usepackage{pifont}
\newcommand{\xmark}{\ding{55}} % Định nghĩa dấu X-mark

\usepackage{parskip}          % Tạo khoảng cách giữa các đoạn văn
\usepackage{pgfgantt}         % Vẽ biểu đồ Gantt

% --- 2. QUẢN LÝ TÀI LIỆU THAM KHẢO ---
\usepackage{biblatex}
\addbibresource{ref.bib}

% --- 3. MÀU SẮC & CODE ---
\usepackage[smartEllipses]{markdown}
\usepackage{color}
\usepackage{fancyvrb}
\usepackage{fvextra}
\usepackage{tcolorbox}
\usepackage{xcolor,colortbl}

% Định nghĩa màu
\definecolor{barblue}{RGB}{153,204,254}
\definecolor{groupblue}{RGB}{51,102,254}
\definecolor{linkred}{RGB}{165,0,33}
\definecolor{dkgreen}{rgb}{0,0.6,0}
\definecolor{xam}{rgb}{0.95, 0.95, 0.95}
\definecolor{codegray}{rgb}{0.5,0.5,0.5}
\definecolor{lightblue}{rgb}{0.8, 0.35, 0.9}
\definecolor{light-gray}{gray}{0.92}  
\definecolor{codegreen}{rgb}{0,0.6,0}
\definecolor{codepurple}{rgb}{0.58,0,0.82}

% Style cho code
\lstdefinestyle{mystyle}{
  language=python,
  backgroundcolor=\color{light-gray}, 
  commentstyle=\color{codegreen},
  keywordstyle=\color{blue},
  numberstyle=\tiny\color{codegray},
  stringstyle=\color{codepurple},
  basicstyle=\ttfamily\footnotesize,
  breakatwhitespace=false,         
  breaklines=true,                 
  captionpos=b,                    
  keepspaces=true,                 
  numbers=left,                    
  numbersep=5pt,                  
  showspaces=false,                
  showstringspaces=false,
  showtabs=false,                  
  tabsize=2
}
\lstset{style=mystyle}

% --- 4. CẤU HÌNH TRANG (GEOMETRY) ---
\usepackage{geometry} 
\usepackage[Bjornstrup]{fncychap}

% Các gói phụ trợ khác
\usepackage[makeroom]{cancel}
\usepackage{nicefrac}
\usepackage{multicol,longtable,amscd}
\usepackage{diagbox}
\usepackage{booktabs}
\usepackage{alltt}
\usepackage{tikz}
\usetikzlibrary{arrows,snakes,backgrounds,calc}
\usepackage{caption}
\usepackage{lastpage}
\usepackage{enumerate}
\usepackage{multirow}
\usepackage{rotating}
\usepackage{setspace}

\usepackage[hidelinks]{hyperref}

% --- 5. ĐỊNH NGHĨA LAYOUT CHÍNH (THESIS LAYOUT) ---
\def\thesislayout{
	\geometry{
		a4paper,
		total={160mm,240mm}, 
		left=30mm,
		top=30mm,
        bottom=25mm,
        right=20mm
	}
}
\thesislayout

% --- 6. CÁC THIẾT LẬP KHÁC ---
\setcounter{secnumdepth}{4}
\setcounter{tocdepth}{3}

\makeatletter
\newcounter {subsubsubsection}[subsubsection]
\renewcommand\thesection{\thechapter.\arabic{section}}
\renewcommand\thesubsubsubsection{\thesubsubsection .\@alph\c@subsubsubsection}
\newcommand\subsubsubsection{\@startsection{subsubsubsection}{4}{\z@}%
                                     {-3.25ex\@plus -1ex \@minus -.2ex}%
                                     {1.5ex \@plus .2ex}%
                                     {\normalfont\normalsize\bfseries}}
\newcommand*\l@subsubsubsection{\@dottedtocline{3}{10.0em}{4.1em}}
\newcommand*{\subsubsubsectionmark}[1]{}
\makeatother

\usepackage{algorithm}
\sloppy
\captionsetup[figure]{labelfont={small,bf},textfont={small},belowskip=-1pt,aboveskip=-9pt}
\captionsetup[table]{labelfont={small,bf},textfont={small},belowskip=-1pt,aboveskip=7pt}

\setlength{\floatsep}{5pt plus 2pt minus 2pt}
\setlength{\textfloatsep}{5pt plus 2pt minus 2pt}
\setlength{\intextsep}{10pt plus 2pt minus 2pt}
\renewcommand{\algorithmicrequire}{\textbf{Input:}}
\renewcommand{\algorithmicensure}{\textbf{Output:}}